\section{Introduction} \label{sec-intro}

The Fiat--Shamir transform~\cite{C:FiaSha86} turns a three-move identification protocol into a signature scheme by computing the verifier's challenge as a hash of the first message and the signed message. Applied to the Schnorr protocol it gives Schnorr signatures~\cite{C:Schnorr89}, among the most widely used pieces of cryptography today, typically in the form of EdDSA~\cite{eddsa}, deployed in TLS and SSH and, over secp256k1, in Bitcoin since the Taproot upgrade~\cite{add:BIP340}. Applied to a lattice protocol with rejection sampling it gives Dilithium~\cite{Dilithium}, standardized by NIST in 2024 as ML-DSA in FIPS~204~\cite{FIPS204}. In both cases the security analyses that support deployment model the hash function as a random oracle~\cite{CCS:BelRog93}, sometimes with further idealizations of the group~\cite{EC:FucPloSeu20,JMC:NevSmaWar09,EPRINT:Shoup23,C:CLMQ21}. In reality the hash function is fixed and public, and its code is available to everyone. We ask what survives when the random oracle is \emph{instantiated}, that is, replaced by an actual hash function. By an instantiation we mean an efficiently evaluable keyed hash family, generated honestly and published in the verification key, whose key may be generated together with the signing key and may depend on the verification key. Our main adaptive Schnorr result uses deterministic nonce generation but retains the usual signature format, and its public hash first applies a TLF to the message; we state both changes where they occur.
\begin{quote}
\emph{Can deployed Fiat--Shamir signature schemes be proved secure in the standard model, and under what assumptions?}
\end{quote}

The question is known to be hard, and the two lines of work that address it leave exactly these schemes untouched. On the negative side, Fiat--Shamir is uninstantiable in general: Goldwasser and Kalai~\cite{FOCS:GolKal03} construct protocols for which the transform is insecure with \emph{every} concrete hash function. On the positive side, correlation-intractable hashing~\cite{JACM:CanGolHal04} instantiates Fiat--Shamir for statistically sound interactive proofs, and such hash functions are now known from standard assumptions~\cite{EC:CCRR18,STOC:CCHLRRW19,C:PeiShi19}. That route does not reach Schnorr or Dilithium. Correlation intractability rules out a prover finding a first message whose challenge is \emph{bad}, meaning one that would let a false statement be accepted. No false statements arise in the relations underlying these two schemes. Every group element has a discrete logarithm, and every honestly generated Dilithium verification key has a short secret. Soundness is therefore vacuous, and there is no bad-challenge relation for correlation intractability to act on. What a signature scheme needs instead is unforgeability, and the random-oracle proofs obtain it by extracting from the forger, either by programming the oracle or by rewinding it. Neither operation is available for a fixed hash function.

So the instantiability of the schemes people actually use is unresolved. For Schnorr, essentially one standard-model result exists: Paillier and Vergnaud~\cite{AC:PaiVer05} showed that for \emph{every} hash function, Schnorr signatures resist \emph{key recovery} under chosen-message attack, a notion they call BK-CMA, assuming the one-more discrete-logarithm (OMDL) problem~\cite{JC:BNPS03} is hard. OMDL is interactive, since the adversary queries an oracle that computes discrete logarithms, so it is not falsifiable in the sense of Naor~\cite{C:Naor03} and resists cryptanalytic validation. There are also meta-reduction barriers. Paillier and Vergnaud rule out standard-model proofs of Schnorr unforgeability from DL-type assumptions for algebraic reductions~\cite{AC:PaiVer05}, and later work rules out tight generic reductions in the ROM~\cite{C:GarBhaLok08,EC:Seurin12,JC:FleJagSch19}. For Dilithium we are not aware of any standard-model security proof. To our knowledge, no standard-model unforgeability proof for Schnorr under an honestly generated concrete hash family was previously known. Prior work gives key recovery for arbitrary hash functions~\cite{AC:PaiVer05}, necessary hash properties with generic-group sufficiency~\cite{JMC:NevSmaWar09}, and positive Fiat--Shamir instantiation for highly sound protocols under a bounded notion~\cite{DBLP:journals/tcs/MittelbachV18}.

This paper addresses the question directly. We give standard-model instantiations under non-interactive assumptions, for Schnorr and for a Lyubashevsky-style lattice scheme motivated by Dilithium. We work with a keyed hash family whose key is generated honestly and published in the verification key; how our reductions pass the meta-reduction barriers is explained at the end of Section~\ref{sec-results}.

\subsection{Our Results} \label{sec-results}

We begin with an adaptive key-recovery result for classical Schnorr. We then turn to unforgeability, first proving bounded static-message security for the classical randomized signer and then adaptive security for Deterministic Schnorr. We carry the static stages to the lattice setting. In each construction the hash function is keyed. The hash key is generated with the signing key and included in the verification key, as in the full-domain-hash instantiation of Hohenberger, Sahai, and Waters~\cite{EC:HohSahWat14}.

\heading{Key recovery and the limit of lossiness.} Our first result concerns BK-CMA, where the adversary receives an adaptive signing oracle but wins only by recovering the signing key. The hash first applies a tunable lossy function to the encoded pair $(R,m)$ and then applies an efficient map to $\Zp$. The real scheme generates the tunable lossy function in injective mode. In the proof, indistinguishability lets the reduction switch to a lossy mode whose image can be enumerated.

The enumerable image gives an exact signing simulation. Let $C\subseteq\Zp$ contain every challenge produced in lossy mode. To sign $m$ without the secret key, the reduction samples $c\getsr C$ and $s\getsr\Zp$, sets $R\gets g^sX^{-c}$, and accepts if $H(R,m)=c$. There is exactly one accepting pair $(c,s)$ for every $R\in\GG$. Thus a trial succeeds with probability $1/|C|$, and conditioned on success the output is distributed exactly as an honest Schnorr signature. The reduction can answer adaptive signing queries and, if the adversary recovers the signing key, solve discrete logarithm (Theorem~\ref{thm-bk}).

The theorem permits any efficient final map to $\Zp$, including a constant map. This is consistent with the security notion: a constant challenge makes Schnorr trivial to forge, but signatures may still reveal no efficient way to recover $x$. One could even omit the mode switch and deploy a fixed highly compressing function whose small range is efficiently known. The same sampler would then prove key-recovery security directly from DL. Such a result would say little about Schnorr as used in practice.

The reason is that a practical challenge hash is intended to spread its outputs throughout the challenge space, not to have a publicly computable polynomial- or quasi-polynomial-size cover. The point of the two-mode construction is that only the lossy mode, used in the proof, has the small enumerable image; the deployed inner function is generated in the injective mode, and the composite has the usual Schnorr domain and range. This does not prove that SHA-3 has a lossy mode or that plain SHA-3 satisfies our theorem. It gives a standard-model family with the same external input and output types without building the small image into the deployed inner function.

The security bound uses both indistinguishability of the two modes and DL. For any fixed image bound, the mode switch need only resist a polynomial-time distinguisher, while the DL reduction pays for enumerating the lossy image and rejection-sampling each signature. A fixed quasi-polynomial or subexponential image bound therefore gives a corresponding continuum of DL assumptions. We give both DDH- and LWE-based TLFs in Section~\ref{sec-lf}. Zhandry's DDH cascade gives the stronger endpoint in which the injective distribution is fixed and the proof chooses a sufficiently large polynomial image bound after fixing the adversary; this yields ordinary DL in the Schnorr group under exponential DDH, or constant-$k$ exponential $k$-linear, in auxiliary groups. The small image also marks the limit of the method. If a small covering set could be efficiently computed from the deployed verification key, the same loop would be a no-query forger running in time about the size of that set. Hidden lossiness therefore gives a key-recovery result, but lossiness alone cannot give unforgeability. This is why the next result uses a different tool.

\heading{Unforgeability from CDL.} Our first unforgeability hash is an obfuscated circuit computing the masked expression of Cho, Fuchsbauer, O'Neill, and Sefranek, so a forgery gives a CDL solution. For declared signing messages, the proof chooses simulated signatures and hardwires their hash values into the circuit. This gives bounded static-message unforgeability while keeping the classical Schnorr signing algorithm and signature format (Theorem~\ref{thm-uf}). The bound is real: the scheme is indexed by the a priori query bound $Q$ because its circuit must be padded to hold $Q$ hardwired entries.

The stronger construction adds a trigger branch. For each message digest $d$, this branch has pseudorandom values $c_d,s_d$ and a point $W_d$ chosen so that $(W_d,s_d)$ is a valid signature; at $W_d$ the circuit returns $c_d$. The reduction can therefore answer a signing query without the secret key, while an honest signer reaches the trigger only with negligible probability. A fixed $W_d$ gives only one signature per message, however, whereas a randomized signer gives fresh signatures when a message is repeated. Our main result therefore uses Deterministic Schnorr: the nonce is a pseudorandom function of the signing key and message, so repeated queries receive the same signature in the real scheme and the simulation (Theorem~\ref{thm-der-adaptive}). The result is adaptive EUF-CMA, with no query bound in the scheme and for messages of any polynomial length. The signature format and Schnorr verification equation are unchanged, while the verification key additionally carries the public hash key used in all of our instantiations.

The proof uses iO to switch the published circuit to functionally equivalent hardwired forms, and puncturable PRFs to replace the selected masks and trigger values by uniform ones. Security rests on the circular discrete-logarithm problem of Cho, Fuchsbauer, O'Neill, and Sefranek~\cite{C:CFOS25}, indistinguishability obfuscation~\cite{C:BGIRSVY01,SIAM:BGIRSVH16,STOC:JaiLinSah21}, and puncturable PRFs~\cite{AC:BonWat13,EC:BoyGilIsh15,CCS:KPTZ13}, together with a TLF applied to the message and a standard-secure PRF for nonce derivation. A fixed quasi-polynomial image bound gives the quantitative hardness profile stated in the theorem. With a strongly efficiently enumerable TLF, the proof may instead choose a sufficiently large polynomial image bound after fixing the adversary, and ordinary PPT security of CDL, iO, and the PRFs suffices. The proof must know a digest before publishing its temporary circuit. Switching the TLF to lossy mode lets it guess this digest from an enumerable set of size $r$, rather than guessing a message from a space of exponential size.

\heading{Lattice signatures.} Our remaining results extend the program to lattice-based Fiat--Shamir with aborts, the paradigm of Lyubashevsky signatures~\cite{AC:Lyubashevsky09,EC:Lyubashevsky12} that Dilithium~\cite{Dilithium} instantiates. We treat an unstructured scheme of that shape, over a sufficiently large prime modulus with prefix-one binary-decomposition challenges. We instantiate the hash function with an obfuscated masked-decomposition circuit and prove bounded static-message unforgeability (Theorem~\ref{thm-dil}). The flooding condition makes rejection negligible, allowing the proof to abort unless the first attempt is used; the remaining signing points are then fixed by presampled nonces before any hash value is computed. With those points fixed, iO hides the switch to a functionally equivalent circuit with their hash values hardwired, puncturable-PRF security replaces the selected masks by uniform values, and LWE replaces the masked inputs by uniform ones; a forgery gives a circular-SIS solution, which a two-line bit-decomposition identity converts to SIS (Lemma~\ref{lem-csis-sis}). For the deterministic signer, repeated queries replay the same response, so one PRF-derived trigger per message suffices. Recomputing that trigger replaces the hardwired table and removes $Q$, giving ordinary EUF-static security with no query bound (Theorem~\ref{thm-dil-der}). The resulting assumption profile is SIS, LWE, iO, and puncturable PRFs, all polynomially secure, at a quasi-polynomial modulus; the modulus condition makes signing rejections negligibly rare, which is what lets the proof fix its programmed points in advance. A standard-model hash function also removes the QROM question that idealized-model analyses of post-quantum schemes must face. Our scheme is not Dilithium. It uses plain matrices rather than the ring structure, prefix-one decomposition challenges rather than the sparse set $B_\tau$, a generic sufficiently large prime rather than ML-DSA's fixed NTT-friendly prime, and no compression or hints. Section~\ref{sec-dilithium} states the comparison in full, and Section~\ref{sec-conclusion} records the gaps as open problems. What we claim transfers to Dilithium is the instantiation technique.

The paper is organized in two parts after the shared preliminaries. Part~I gives the three Schnorr results in order: adaptive key recovery for classical Schnorr, bounded static-message unforgeability for classical Schnorr, and adaptive unforgeability for Deterministic Schnorr. The appendix proves the adaptive deterministic result. Part~II gives the lattice preliminaries and scheme, bounded static security, static security without a query bound for deterministic signing, and the path to sparse challenges and adaptive security. Section~\ref{sec-conclusion} records the remaining problems.

Our reductions are non-black-box in the hash function: they use the code of the PRF keys and, in the unforgeability constructions, choose the hash function's circuit. This is how the results evade the meta-reduction barriers of~\cite{AC:PaiVer05,C:GarBhaLok08,EC:Seurin12,JC:FleJagSch19}, which apply to reductions treating the assumption in a representation-independent way, and the non-programmable-ROM impossibilities of Fukumitsu and Hasegawa~\cite{FukumitsuH21}, which simulate forgers that make random-oracle queries. CDL is representation-dependent by design, and our forgers make no oracle queries.

\subsection{Technical Overview} \label{sec-techov}

\heading{The classical proof, and what breaks.} Recall Schnorr signatures over a group $\GG = \langle g \rangle$ of prime order $p$. The secret key is $x \getsr \Zp$ with public key $X = g^x$; a signature on $m$ is $(R, s)$ with $R = g^\rho$, $c = H(R, m)$, and $s = \rho + cx \bmod p$; verification checks $g^s = R \cdot X^c$. In the ROM proof, the reduction answers a signing query by sampling $s, c \getsr \Zp$, setting $R \gets g^s X^{-c}$, and \emph{programming} the random oracle so that $H(R,m) = c$. A standard-model hash function cannot be programmed, as it is fixed before the adversary runs. Our proofs instead arrange that the simulated signatures already verify under the fixed function.

\heading{An obfuscated hash.} For unforgeability, the hash value in a forgery must give a CDL solution. We start from the ROM proof of~\cite{C:CFOS25}. Its reduction answers a hash query $(R,m)$ with $f(Rh^ag^b)+a$ for fresh $a,b\getsr\Zp$, where $h$ is the CDL challenge and $f$ is the conversion function. These values are uniform, and the known pair $(a,b)$ converts a forgery to a CDL solution. Our hash function is an obfuscated circuit that computes the same expression, with $X$ in place of $h$ and with $a,b$ computed by puncturable PRFs on $(R,m)$.

Signing queries are handled by hardwiring. Since the messages are declared before the verification key is published, the proof samples $s_i,c_i\getsr\Zp$, sets $R_i=g^{s_i}X^{-c_i}$, and hardwires $(R_i,m_i)\mapsto c_i$. The iO step hides the switch to this functionally equivalent hardwired circuit; puncturable-PRF security then lets the proof replace the PRF outputs at the hardwired points by uniform values. Before applying this change of variables, the proof replaces the honest hash value at $(R_i,m_i)$ by an independent uniform value. This replacement is exact because the two PRF masks at that point are uniform and occur nowhere else. A forgery is on a fresh message, so it is evaluated by the PRF-computed part of the circuit and gives a CDL solution. We choose $f$ to have no zeros, ensuring that the CDL game accepts the resulting solution. The hardwired entries make the proof circuit grow linearly with an a priori signing-query bound $Q$ built into this basic scheme.

\heading{A trigger for signing.} The deterministic construction places a trigger before the main branch. For a message $m$, let $d=\dig(K,m)$ be its digest under the public TLF key. Two PRFs define
$$
c_d=\PRFEv_3(k_3,\enczero{d}),
\qquad
s_d=\PRFEv_4(k_4,\enczero{d}),
\qquad
W_d=g^{s_d}X^{-c_d}.
$$
On input $(R,d)$ the circuit returns $c_d$ if $R=W_d$, and otherwise evaluates the main branch above. Thus
$$
W_d X^{H(W_d,m)}=W_dX^{c_d}=g^{s_d},
$$
so $(W_d,s_d)$ is a valid signature on $m$ and the reduction can sign without knowing $x$. Pseudorandomness of the nonce PRF first lets the proof replace the honest nonce by a uniform value. In that hybrid the honest signing point equals $W_d$ with probability $1/p$, so the trigger affects honest use only negligibly. In the signing hybrid, the joint distribution of the honest and trigger signatures is made symmetric, and the two are exchanged. A main-branch forgery gives a CDL solution using the full main-branch PRF keys. For a trigger-branch forgery, the reduction guesses its digest, keeps $\PRF_3$ so that the trigger challenge remains known, and punctures only the response PRF $\PRF_4$ to embed a DL challenge.

\heading{Repeated queries and deterministic signing.} One trigger point gives one signature per digest. If a randomized signer is queried twice on the same message, the real scheme returns two fresh signatures, while the reduction would return $(W_d,s_d)$ twice. The main construction removes this mismatch by deriving the nonce from a PRF of the signing key and message. The real signer is then deterministic as well, so a repeated query is answered by replaying the same signature. For any stateless deterministic signature scheme, this replay argument generically reduces full EUF-CMA security to adaptive security against adversaries that never repeat a message. It does not reduce our theorem to the basic classical Schnorr result, however, because that result fixes all messages before key generation and builds a polynomial $Q$ into the scheme. The trigger branch and local hybrids establish security for an adaptive sequence of distinct messages without such a scheme-level parameter, while determinism handles only its repetitions.

\heading{The local signing hybrid.} For adaptive security, the hash first applies a TLF to the message. Let $K$ be its public key, let $\dig(K,m)$ denote its encoded output on $m$, let $m_j$ be the $j$-th distinct signing message, and write
$$
d_j=\dig(K,m_j).
$$
The digest $d_j$ is the input at which the trigger PRFs derive $c_{d_j},s_{d_j},W_{d_j}$. Let $(R_j,c_j,s_{0,j})$ be the honest record, where $c_j$ is the hash challenge and $s_{0,j}$ the response, and let $(W_{d_j},c_{d_j},s_{d_j})$ be the trigger record. Define $\Hm_i$ to be the complete experiment in which the first $i$ distinct signing messages receive trigger signatures and the rest receive honest signatures. Every $\Hm_i$ publishes the original, unpunctured circuit.

Fix $j$ and condition on the proof having correctly guessed $d_j$. The proof first excludes the event $R_j=W_{d_j}$, which has probability $1/p$ and is also what makes the circuit switch below functionally valid. To compare $\Hm_{j-1}$ and $\Hm_j$, iO switches the published circuit to a functionally equal temporary circuit carrying exactly the encodings of the following two hash input--output pairs:
$$
\big((R_j,d_j),c_j\big)
\qquad\text{and}\qquad
\big((W_{d_j},d_j),c_{d_j}\big).
$$
The two main-branch PRFs are punctured at the encodings of these two inputs, and the two trigger PRFs are punctured at the one input encoding $d_j$. The main-branch punctures at $R_j$ hide the masks used to compute the honest challenge, letting the proof make that challenge uniform. The main-branch PRFs are also punctured at $W_{d_j}$ even though the main branch is not evaluated there; this keeps the puncture set from labeling $R_j$ as the honest point. The trigger-branch punctures hide the trigger challenge and response, letting the proof make both uniform. The two explicit pairs preserve the circuit's output at $R_j$ and $W_{d_j}$ while the punctured keys evaluate it everywhere else. Puncturable-PRF security and exact changes of variables then turn the records into
$$
\big(g^{s_{0,j}}X^{-c_j},c_j,s_{0,j}\big)
\qquad\text{and}\qquad
\big(g^{s_{d_j}}X^{-c_{d_j}},c_{d_j},s_{d_j}\big).
$$
Before conditioning on distinct first components, the four scalars are uniform and independent and the two records are generated by the same rule. The condition $R_j\ne W_{d_j}$ is itself invariant under exchange, so the conditioned joint distribution remains invariant under exchange, although the records are then correlated. The temporary circuit stores their two input--output pairs in a canonical unordered set, and the proof discards every other copy or label that distinguishes them. It may therefore exchange which response it returns, and then reverse the changes. The endpoint is $\Hm_j$, again with the original circuit.

\heading{Why there is no query bound in the scheme.} The temporary circuit for the $j$-th distinct message disappears before the next distinct message is treated. The next comparison starts again from the original circuit and uses a new temporary circuit with only two pairs. Thus punctures and hardwired values never accumulate. This local argument removes the a priori parameter $Q$ from the scheme. The reduction still runs one hybrid per distinct query, so its loss depends on the adversary's polynomial number $q_s$ of signing queries, but the scheme and circuit do not.

\heading{What lossiness does.} The local hybrid must fix the digest $d_j$ before publishing its temporary circuit. After replacing the nonce PRF by a random function, the proof can presample the uniform nonce for the $j$-th new message, but it still does not know the adaptive message $m_j$ or its digest. Switching the TLF to lossy mode lets the reduction guess $d_j$ from an enumerated image of size at most $r$, rather than guess $m_j$ at cost $2^{\msglen}$. Lossiness therefore handles adaptivity; it is not what removes $Q$.

\subsection{Related Work} \label{sec-relwork}

\heading{Schnorr in idealized models.} Pointcheval and Stern~\cite{EC:PoiSte96,JC:PoiSte00} proved EUF-CMA security of Schnorr signatures in the ROM under DL, with a loose reduction. Tight reductions are known in the ROM combined with the algebraic group model~\cite{EC:FucPloSeu20} and in generic group models~\cite{JMC:NevSmaWar09,EPRINT:Shoup23,C:CLMQ21}. In the ROM alone, reductions are known under the interactive MBDL assumption~\cite{INDOCRYPT:BelDai20} and, with some loss, under DDH-type assumptions~\cite{C:RotSeg21}.

CFOS~\cite{C:CFOS25} give a completely tight ROM reduction under CDL, a non-interactive and falsifiable assumption. Our Schnorr construction builds directly on their simulation strategy.

\heading{Barriers.} Meta-reductions rule out tight generic reductions for Schnorr in the ROM~\cite{AC:PaiVer05,C:GarBhaLok08,EC:Seurin12,JC:FleJagSch19}, and rule out standard-model proofs in the non-programmable ROM for certain reduction classes~\cite{FukumitsuH21}. Both classes cover reductions that treat the assumption in a representation-independent way and run the forger as a black box. Our reductions do neither, as explained at the end of Section~\ref{sec-results}. Mittelbach and Venturi~\cite{DBLP:journals/tcs/MittelbachV18} give positive standard-model Fiat--Shamir results for highly sound protocols, using obfuscation and puncturable PRFs and a query-bounded security notion; their setting is the soundness one described above, and the bounded notion is the same shape as our Corollary~\ref{cor-uf}. Neven, Smart, and Warinschi~\cite{JMC:NevSmaWar09} identify standard-model properties that a Schnorr hash function must have, such as random-prefix preimage resistance, and show sufficiency in the generic group model; they do not give a standard-model unforgeability proof. The uninstantiability results for Fiat--Shamir~\cite{FOCS:GolKal03}, and for random-oracle schemes generally~\cite{JACM:CanGolHal04}, use contrived constructions, and the correlation-intractability line~\cite{EC:CCRR18,STOC:CCHLRRW19,C:PeiShi19} covers statistically sound proofs; Section~\ref{sec-intro} explains why neither reaches the schemes treated here.

\heading{Instantiating random oracles.} Hohenberger, Sahai, and Waters~\cite{EC:HohSahWat14} instantiate full-domain hash for RSA and BLS signatures using iO and puncturable PRFs. Their selective result fixes only the forgery message before key generation and allows adaptive signing queries, whereas our static result fixes the signing messages as well. Both proofs puncture PRFs, but their signing simulations differ. There are two further differences. For Schnorr, the signer chooses the group element $R$ in each signature, and the proof uses this freedom by choosing those points in advance and hardwiring their hash values, in the manner of honest-verifier zero-knowledge simulation. Also, the forgery is converted to a CDL solution rather than inverted using the trapdoor structure of RSA/BLS. Lossy functions were introduced by Peikert and Waters~\cite{STOC:PeiWat08,JACM:PeiWat11}. Zhandry~\cite{C:Zhandry16} introduced their extremely lossy variant and used it to instantiate several random-oracle properties; Fischlin and Rohrbach~\cite{TCC:FisRoh23} showed ELFs cannot be built from a random oracle. Murphy, O'Neill, and Zaheri~\cite{AC:MurONeZah22} gave instantiation results for OAEP-type and Fujisaki--Okamoto-type transforms using obfuscation-based tools. We are not aware of prior standard-model instantiations of Schnorr signatures beyond~\cite{AC:PaiVer05}.

\heading{Keyed hashing in the verification key.} In all our constructions the hash function used for the instantiation is keyed. The BK-CMA hash generator ignores the Schnorr public key, so its hash key is independent of the signing key, although each user still receives a separately generated public hash key. In our unforgeability constructions the hash key instead depends on the public key of the signer, which is embedded in an obfuscated program. This is undesirable in that we would like to show an independent hash key works, which is more consistent with cryptographic practice (e.g., SHA is not designed with Schnorr in mind). However, at present this seems technically difficult.

Nevertheless, we stress that such instantiations \emph{still refute uninstantiability}. For example the CGH theorem is stated for an independently sampled hash key, but the proof extends to public-key-dependent or jointly generated instantiations, since the concrete attack uses only public evaluation of the hash function, not independence from other keys of the overlying scheme. Moreover, such a restriction was also present in prior instantiability work~\cite{EC:HohSahWat14}.
